\section{Smooth orbit averages and the pointwise theorem}
\label{sec:ergodic}

This section deduces Theorem~\ref{thm:main} from a finite oscillation
estimate for smooth orbit averages.  The oscillation estimate forces
almost everywhere convergence on every invertible probability preserving
system.  The mixing hypothesis \eqref{eq:mixing} identifies the limit, and
signed approximations then recover all positive integer averaging lengths.

We use the notion of an invertible probability preserving system fixed in
the introduction.  Every integer power of \(T\) is measurable and measure
preserving: positive powers follow by iteration, and applying invariance to
the measurable set \(TA\) gives \(\mu(TA)=\mu(A)\).

\begin{definition}[Averaging profiles]
\label{def:averaging-profile}
An averaging profile is a real function
\(\varphi\in C_c^\infty(\mathbb R;\mathbb R)\) such that, for some
\(G\geq1\),
\begin{equation}
 \|\varphi^{(n)}\|_\infty
   \leq\bigl(G(n+2)\bigr)^{G(n+2)}
       \qquad(n=0,1,2,\ldots),
 \label{eq:profile-derivatives}
\end{equation}
and
\begin{equation}
 \int_{\mathbb R}\varphi(t)\,dt=1,\qquad
 \int_{\mathbb R}t^j\varphi(t)\,dt=0\quad(1\leq j\leq3).
 \label{eq:profile-moments}
\end{equation}
The function may have either sign.  For \(a>0\), write
\(\varphi_a(t)=a^{-1}\varphi(t/a)\).
\end{definition}

Section~\ref{sec:transfer} uses the moment and derivative conditions to
construct local kernels for adjacent smooth scale differences.
Lemma~\ref{lem:ergodic-moments} shows that these restrictions still permit
signed profiles arbitrarily close in \(L^1\) to
\(h_c(t)=c^{-1}\mathbf1_{(0,c]}(t)\) for every \(c\in[1,2]\).

\begin{proposition}[Finite oscillation of smooth orbit averages]
\label{prop:orbit-oscillation}
For every averaging profile \(\varphi\), there is a constant
\(C_\varphi<\infty\) with the following property.  On every invertible
probability preserving system, let \(f_1,f_2,f_3:X\to\mathbb C\) be
measurable with \(|f_j(x)|\leq1\) for all \(x\), and define
\[
 V_k^\varphi(x)=\sum_{n\in\mathbb Z}\varphi_{2^k}(n)
                 \prod_{j=1}^3f_j(T^{jn}x),
       \qquad k=0,1,2,\ldots .
\]
For every integer \(w\geq1\) and every deterministic list of integers
\(0\leq k_1<\cdots<k_{w+1}\),
\begin{equation}
 \sum_{i=1}^w\int_X
   \max_{k_i\leq\ell\leq k_{i+1}}
      |V_\ell^\varphi-V_{k_i}^\varphi|\,d\mu
       \leq C_\varphi\sqrt w.
 \label{eq:orbit-oscillation}
\end{equation}
The constant is independent of the system, the inputs, \(w\), and the
endpoints.
\end{proposition}

Section~\ref{sec:transfer} proves Proposition~\ref{prop:orbit-oscillation}
from Theorem~\ref{thm:depth-local}, which the later analytic sections prove.
The sums defining \(V_k^\varphi\) are finite.  If
\(\operatorname{supp}\varphi\subseteq[-B,B]\) with \(B\geq1\), then
\begin{equation}
 |V_k^\varphi(x)|
 \leq 2^{-k}(2B2^k+1)\|\varphi\|_\infty
 \leq(2B+1)\|\varphi\|_\infty.
 \label{eq:smooth-uniform-bound}
\end{equation}
Thus the averages and all the finite maxima in
\eqref{eq:orbit-oscillation} are measurable and uniformly bounded.

The next criterion is the finite-block oscillation argument used by
Bourgain and Demeter \cite{Bourgain1989,Demeter2007}.
\begin{lemma}[Oscillation forces almost everywhere convergence]
\label{lem:ergodic-cauchy}
Let \((X,\mathcal F,\mu)\) be any probability space and let
\(Z_k:X\to\mathbb C\), \(k\geq0\), be measurable.  Suppose that
\(|Z_k(x)|\leq M<\infty\) for every \(k,x\), and that there is
\(C<\infty\) such that, for every integer \(w\geq1\) and every
deterministic list \(0\leq k_1<\cdots<k_{w+1}\),
\[
 \sum_{i=1}^w\int_X
   \max_{k_i\leq\ell\leq k_{i+1}}|Z_\ell-Z_{k_i}|\,d\mu
       \leq C\sqrt w.
\]
Then \(Z_k(x)\) converges in \(\mathbb C\) on a measurable conull set.
\end{lemma}

\begin{proof}
Define
\[
 \omega(x)=\lim_{K\to\infty}\sup_{p,q\geq K}|Z_p(x)-Z_q(x)|.
\]
The suprema are countable and decrease with \(K\), so \(\omega\) is
measurable.  Completeness of \(\mathbb C\) says that the convergence set
is exactly \(\{\omega=0\}\).  If its complement has positive measure,
then, since
\(\{\omega>0\}=\bigcup_{h\geq1}\{\omega>2/h\}\), there are
\(\delta>0\) and a measurable set
\(E=\{\omega>2\delta\}\) of measure \(d>0\).  For every \(K\geq0\) and
\(x\in E\),
\[
 \sup_{\ell\geq K}|Z_\ell(x)-Z_K(x)|
 \geq\frac12\sup_{p,q\geq K}|Z_p(x)-Z_q(x)|
 \geq\frac12\omega(x)>\delta.
\]
The integral of the first supremum is at least \(\delta d\).
For \(L\geq K\), the measurable finite maxima
\[
 M_{K,L}(x)=\max_{K\leq\ell\leq L}|Z_\ell(x)-Z_K(x)|
\]
increase to this supremum.  By monotone convergence there is a
deterministic \(L>K\) with \(\int_XM_{K,L}\,d\mu\geq\delta d/2\).
Starting with \(k_1=0\), choose \(k_{i+1}>k_i\) successively with this
property.  The assumed estimate gives
\(w\delta d/2\leq C\sqrt w\) for every \(w\), a contradiction for large
\(w\).  Hence \(\omega=0\) almost everywhere.
\end{proof}

Proposition~\ref{prop:orbit-oscillation} and
Lemma~\ref{lem:ergodic-cauchy} already give an almost everywhere limit
for each fixed averaging profile.  We now identify it for mixing systems.

\subsection{The product limit in norm}

Suppose for this subsection that \(T\) satisfies the mixing condition
\eqref{eq:mixing}.
For every fixed bounded measurable \(g,h:X\to\mathbb C\), this implies
\begin{equation}
 \int_X g(x)h(T^m x)\,d\mu(x)
   \longrightarrow
      \left(\int_Xg\,d\mu\right)\left(\int_Xh\,d\mu\right)
       \qquad(|m|\to\infty).
 \label{eq:ergodic-function-mixing}
\end{equation}
For finite complex simple functions this is the finite expansion into
the set correlations in \eqref{eq:mixing}.  Every bounded
complex measurable function has uniform simple-function approximations,
obtained by partitioning its bounded range into finitely many small sets.
If \(g_0,h_0\) are such approximations, the correlation error is at most
\[
 \|g-g_0\|_\infty\|h\|_\infty+
        \|g_0\|_\infty\|h-h_0\|_\infty
\]
uniformly in \(m\), and the product of integrals has the same continuity
property.  Taking the mixing limit for fixed approximations and then
letting their errors tend to zero proves
\eqref{eq:ergodic-function-mixing}.

We include the classical van der Corput induction for multiple averages,
specialized to mixing; see \cite[Section~3]{FKO1982} and
\cite[Section~3]{Bergelson1987}.  The finite Hilbert space van der Corput
inequality used below is the one recorded in \cite[Lemma~D.1]{HostKra2005};
the criterion follows by taking limits.
\begin{lemma}[Hilbert space van der Corput criterion]
\label{lem:ergodic-vdc}
Let \(\mathcal H\) be a complex Hilbert space, with inner product linear
in the first variable.  Let \(u_n\in\mathcal H\), \(n\geq1\), satisfy
\(\|u_n\|\leq M<\infty\).  Suppose that, for each fixed integer \(h\geq1\),
\[
 \gamma_h=\lim_{N\to\infty}
       \frac1N\sum_{n=1}^{N-h}\langle u_{n+h},u_n\rangle
\]
exists, with a sum over an empty range interpreted as zero.  If
\(\gamma_h\to0\) as \(h\to\infty\), then
\(\|N^{-1}\sum_{n=1}^Nu_n\|\to0\).
\end{lemma}

\begin{proof}
For fixed integers \(N,H\geq1\), extend \(u_n\) by zero outside
\(\{1,\ldots,N\}\).  The identity
\[
 H\sum_{n=1}^Nu_n
   =\sum_{m=1}^{N+H-1}\sum_{j=0}^{H-1}u_{m-j}
\]
and Cauchy's inequality give
\begin{equation}
\begin{split}
 \left\|\frac1N\sum_{n=1}^Nu_n\right\|^2
 \leq\frac{N+H-1}{N^2H^2}\bigg[
 &H\sum_{n=1}^N\|u_n\|^2\\
 &+2\sum_{h=1}^{H-1}(H-h)
       \operatorname{Re}\sum_{n=1}^{N-h}
                \langle u_{n+h},u_n\rangle\bigg].
\end{split}
\label{eq:ergodic-finite-vdc}
\end{equation}
Indeed, after Cauchy's inequality the expression to expand is
\((N+H-1)\sum_m\|\sum_{j=0}^{H-1}u_{m-j}\|^2\).  Each diagonal term
occurs \(H\) times, and each pair of positions at difference \(h\)
occurs \(H-h\) times, giving the displayed bracket.

Fix \(H\) and let \(N\to\infty\) in
\eqref{eq:ergodic-finite-vdc}.  Then
\begin{align*}
 \limsup_{N\to\infty}\left\|\frac1N\sum_{n=1}^Nu_n\right\|^2
 &\leq\frac{M^2}{H}
      +\frac2{H^2}\sum_{h=1}^{H-1}(H-h)\operatorname{Re}\gamma_h\\
 &\leq\frac{M^2}{H}+\frac2H\sum_{h=1}^{H-1}|\gamma_h|.
\end{align*}
The right side tends to zero as \(H\to\infty\), because
\(|\gamma_h|\to0\) implies convergence of its Ces\`aro means to zero.
\end{proof}

\begin{lemma}[Ces\`aro norm convergence at distinct slopes]
\label{lem:ergodic-cesaro}
On every invertible probability preserving system satisfying
\eqref{eq:mixing}, for every integer \(r\geq1\), every
list of pairwise distinct nonzero integers
\(\lambda_1,\ldots,\lambda_r\), and every list of bounded measurable
\(f_1,\ldots,f_r:X\to\mathbb C\),
\[
 \left\|
 \frac1N\sum_{n=1}^N\prod_{j=1}^r f_j(T^{\lambda_jn}\,\cdot)
       -\prod_{j=1}^r\int_Xf_j\,d\mu
 \right\|_{L^2(\mu)}
 \longrightarrow0\qquad(N\to\infty).
\]
\end{lemma}

\begin{proof}
Write \(U^m f=f\circ T^m\), a unitary operator on \(L^2(\mu;\mathbb C)\).
Use the inner product \(\langle f,g\rangle=\int_Xf\overline g\,d\mu\).
We induct on \(r\), taking the empty product to be \(1\), so the
assertion for \(r=0\) is immediate.  Suppose it is known for \(r-1\).
First assume that at least one \(f_j\) has integral zero.  Put
\[
 u_n=\prod_{j=1}^rU^{\lambda_jn}f_j,\qquad
 g_{j,h}=(U^{\lambda_jh}f_j)\overline{f_j},\qquad n,h\geq1.
\]
The vectors \(u_n\) are uniformly bounded in \(L^2\).  Invariance gives
\begin{align*}
 \langle u_{n+h},u_n\rangle
 &=\int_X\prod_{j=1}^rU^{\lambda_jn}g_{j,h}\,d\mu\\
 &=\int_Xg_{1,h}\prod_{j=2}^r
             U^{(\lambda_j-\lambda_1)n}g_{j,h}\,d\mu .
\end{align*}
The second equality shifts all times by \(-\lambda_1n\).
For fixed \(h\), the remaining slopes
\(\lambda_j-\lambda_1\), \(2\leq j\leq r\), are pairwise distinct and
nonzero, and the functions \(g_{j,h}\) are fixed and bounded.  Induction
therefore gives
\[
 \frac1{N-h}\sum_{n=1}^{N-h}
       \prod_{j=2}^rU^{(\lambda_j-\lambda_1)n}g_{j,h}
 \longrightarrow\prod_{j=2}^r\int_Xg_{j,h}\,d\mu
 \quad\text{in }L^2(\mu),\qquad N\to\infty,\ N>h.
\]
For \(r=1\) both products are \(1\), so this assertion is immediate.
The functional \(A\mapsto\int_Xg_{1,h}A\,d\mu\) is continuous on \(L^2\)
by Cauchy's inequality.  Hence, using \((N-h)/N\to1\),
\[
 \lim_{N\to\infty}\frac1N\sum_{n=1}^{N-h}
       \langle u_{n+h},u_n\rangle
   =\prod_{j=1}^r c_j(h),\qquad
 c_j(h)=\int_X(U^{\lambda_jh}f_j)\overline{f_j}\,d\mu .
\]
By \eqref{eq:ergodic-function-mixing} and \(\lambda_j\neq0\),
\[
 c_j(h)\longrightarrow\left|\int_Xf_j\,d\mu\right|^2
       \qquad(h\to\infty).
\]
The factors are bounded by \(\|f_j\|_\infty^2\), and at least one limit
is zero.  Their product tends to zero, so
Lemma~\ref{lem:ergodic-vdc} gives
\(\|N^{-1}\sum_{n=1}^Nu_n\|_2\to0\).

For a general list, let \(m_j=\int_Xf_j\,d\mu\).  The identity
\[
 \prod_{j=1}^rU^{\lambda_jn}f_j-\prod_{j=1}^rm_j
 =\sum_{i=1}^r
       \left(\prod_{j<i}U^{\lambda_jn}f_j\right)
       U^{\lambda_in}(f_i-m_i)
       \left(\prod_{j>i}m_j\right)
\]
expresses the difference as a finite sum of products with a centered
factor.  Regard the later constants as constant functions at their
corresponding slopes.  The centered case just proved applies to every
summand and completes the induction.

For each fixed \(h\), induction was used only for the fixed functions
\(g_{j,h}\).  The van der Corput argument fixes its finite lag range
\(H\) before sending \(N\) to infinity, and sends \(H\) to infinity
afterwards.  No convergence rate uniform in \(h\) is needed.
\end{proof}

\begin{lemma}[Norm convergence with smooth weights]
\label{lem:ergodic-smooth-norm}
On every invertible probability preserving system satisfying
\eqref{eq:mixing}, fix an integer
\(r\geq1\), pairwise distinct nonzero integers
\(\lambda_1,\ldots,\lambda_r\), and bounded measurable
\(f_1,\ldots,f_r:X\to\mathbb C\).  Set
\[
 F_n=\prod_{j=1}^r f_j(T^{\lambda_jn}\,\cdot),\qquad
 p=\prod_{j=1}^r\int_Xf_j\,d\mu,\qquad n\in\mathbb Z.
\]
For every \(\varphi\in C_c^1(\mathbb R;\mathbb C)\) with
\(\int_{\mathbb R}\varphi=1\),
\[
 \left\|\sum_{n\in\mathbb Z}a^{-1}\varphi(n/a)F_n-p\right\|_2
       \longrightarrow0\qquad(a\to\infty,\ a\geq1).
\]
In particular the conclusion holds for every averaging profile.
\end{lemma}

\begin{proof}
Apply Lemma~\ref{lem:ergodic-cesaro} to the slopes \(\lambda_j\) and
then to the slopes \(-\lambda_j\).  For the \(L^2\) vectors
\[
 R_N^\pm=\sum_{n=1}^N(F_{\pm n}-p)
\]
this gives \(\|R_N^\pm\|_2/N\to0\) for both signs.
Choose \(B\geq1\) with
\(\operatorname{supp}\varphi\subseteq[-B,B]\).  For \(a\geq1\), put
\(M_a=\lfloor Ba\rfloor+1\) and
\(w_n^\pm=a^{-1}\varphi(\pm n/a)\), so \(w_{M_a}^\pm=0\).
Vector-valued summation by parts gives
\[
 \sum_{n=1}^{M_a}w_n^\pm(F_{\pm n}-p)
  =\sum_{n=1}^{M_a-1}(w_n^\pm-w_{n+1}^\pm)R_n^\pm.
\]
Since \(|w_n^\pm-w_{n+1}^\pm|\leq a^{-2}\|\varphi'\|_\infty\), the norm
of this sum is at most
\[
 \frac{\|\varphi'\|_\infty}{a^2}
           \sum_{n=1}^{M_a-1}\|R_n^\pm\|_2.
\]
For any \(\epsilon>0\), choose \(N_0\) with
\(\|R_n^\pm\|_2\leq\epsilon n\) for both signs and every \(n\geq N_0\).
As \(M_a/a\leq B+1\),
\[
 \frac1{a^2}\sum_{n=1}^{M_a-1}\|R_n^\pm\|_2
 \leq\frac1{a^2}\sum_{n=1}^{N_0-1}\|R_n^\pm\|_2
       +\frac{\epsilon}{a^2}\sum_{n=1}^{M_a-1}n
 \leq o(1)+\frac{\epsilon}{2}(B+1)^2.
\]
Letting \(a\to\infty\) and then \(\epsilon\downarrow0\) proves convergence
to zero for the positive-time and negative-time parts of the weighted sum
of \(F_n-p\).
The term at \(n=0\) has norm at most
\(a^{-1}|\varphi(0)|\|F_0-p\|_2\), which also tends to zero.
Finally the Riemann sums
\(S_a=a^{-1}\sum_n\varphi(n/a)\) tend to \(\int\varphi=1\), and
\[
 \sum_n a^{-1}\varphi(n/a)F_n-p
   =\sum_n a^{-1}\varphi(n/a)(F_n-p)+p(S_a-1).
\]
This proves the assertion for signed and complex weights alike.
\end{proof}

\subsection{Moment approximation and all lengths}

We need averaging profiles arbitrarily close in \(L^1\) to normalized
interval indicators.  We first fix the smooth bump used in their
construction and later in the local kernels.

\begin{lemma}[A compact Gevrey bump]
\label{lem:ergodic-bump}
There is a real nonnegative \(\rho\in C_c^\infty(\mathbb R)\), supported
in \([-1,1]\), with integral \(1\), whose derivatives obey a bound of
the form \eqref{eq:profile-derivatives} for some constant \(G\geq1\).
\end{lemma}

\begin{proof}
Start with
\[
 \rho_0(x)=
 \begin{cases}e^{-1/(1-x^2)},&|x|<1,\\0,&|x|\geq1.\end{cases}
\]
For \(|x|<1\), put \(d=1-x^2\).  If \(|z-x|\leq c d\) in the complex
plane, with a sufficiently small fixed \(c>0\), then
\(1-z^2=d(1-\theta)\) with \(|\theta|\leq1/4\).  Consequently
\(\operatorname{Re}(1/(1-z^2))\geq c'/d\) for a fixed \(c'>0\).
Cauchy's estimate gives
\[
 |\rho_0^{(n)}(x)|\leq n!(cd)^{-n}e^{-c'/d}.
\]
For every fixed \(n\), the right side tends to zero as \(x\) approaches
either endpoint, so the derivatives extend there by zero and the
extension is smooth.  Maximizing \(d^{-n}e^{-c'/d}\) for \(0<d\leq1\)
and using \(n!\leq(n+1)^n\) gives a bound of the form
\eqref{eq:profile-derivatives}.  Divide by the positive integral of
\(\rho_0\) to obtain \(\rho\).
\end{proof}

\begin{lemma}[Approximation with three vanishing moments]
\label{lem:ergodic-moments}
For every \(c\in[1,2]\) and every \(\epsilon>0\), there is an averaging
profile \(\varphi\) with
\[
 \|\varphi-h_c\|_{L^1(\mathbb R)}<\epsilon,\qquad
 h_c(t)=c^{-1}\mathbf 1_{(0,c]}(t).
\]
The support and derivative constant may depend on \(c,\epsilon\).
\end{lemma}

\begin{proof}
Use the bump \(\rho\) from Lemma~\ref{lem:ergodic-bump}.  For
\(\delta>0\), set \(\rho_\delta(t)=\delta^{-1}\rho(t/\delta)\) and
\(\psi=h_c*\rho_\delta\).  It is real, compactly supported, smooth,
and has mass \(1\).  Its derivatives satisfy
\[
 \|\psi^{(n)}\|_\infty
   \leq\|h_c\|_1\|\rho_\delta^{(n)}\|_\infty
   =\delta^{-n-1}\|\rho^{(n)}\|_\infty,
\]
which has the form \eqref{eq:profile-derivatives} for fixed \(\delta\).
Moreover,
\[
 \|\psi-h_c\|_1
 \leq\int\rho_\delta(s)\|h_c(\,\cdot-s)-h_c\|_1\,ds
 \leq\frac{2\delta}{c}\int|v|\rho(v)\,dv,
\]
because \(\|h_c(\,\cdot-s)-h_c\|_1=2\min(|s|,c)/c\).
Choose \(\delta\) to make this error less than \(\epsilon/2\), fix
\(\psi\), and let \(m_j=\int t^j\psi(t)\,dt\), \(1\leq j\leq3\).

We correct these moments with broad translated bumps: their \(j\)th moments
will be \(D^j\) times fixed numbers, allowing a linear system with a fixed
coefficient matrix to cancel the three moments with total coefficient size
tending to zero as \(D\to\infty\).
Fix distinct real numbers \(z_0,z_1,z_2,z_3\), and for \(D\geq1\) put
\[
 \rho_{i,D}(t)=D^{-1}\rho(t/D-z_i),\qquad
 P_j(z)=\int(z+v)^j\rho(v)\,dv,\qquad 0\leq i,j\leq3.
\]
Then \(\int t^j\rho_{i,D}(t)\,dt=D^jP_j(z_i)\).
Each \(P_j\) is monic of degree \(j\).  The matrix
\(\mathsf M=(P_j(z_i))_{j,i=0}^3\) is therefore the Vandermonde matrix
\((z_i^j)_{j,i}\) multiplied on the left by a lower triangular matrix
with diagonal entries \(1\).  Its determinant is
\(\prod_{i<i'}(z_{i'}-z_i)\neq0\).
Choose the real vector \(\beta=(\beta_0,\ldots,\beta_3)\) satisfying
\[
 \mathsf M\beta=(0,-m_1/D,-m_2/D^2,-m_3/D^3)^{\mathsf T}.
\]
The inverse matrix is fixed, so
\[
 \sum_i|\beta_i|
 \leq\|\mathsf M^{-1}\|_{\ell^1\to\ell^1}
      \left(\frac{|m_1|}{D}+\frac{|m_2|}{D^2}+\frac{|m_3|}{D^3}\right)
 \leq\frac{C_\psi}{D}.
\]
Thus \(\varphi=\psi+\sum_i\beta_i\rho_{i,D}\) has mass \(1\) and its
first three moments vanish exactly.  The extra \(L^1\) error is at most
\(\|\rho\|_1\sum_i|\beta_i|\leq C_\psi/D\), which is less than
\(\epsilon/2\) for large \(D\).  Each \(\rho_{i,D}\) is a fixed translate
and dilation of \(\rho\), with
\(\|\rho_{i,D}^{(n)}\|_\infty=D^{-n-1}\|\rho^{(n)}\|_\infty\).
Multiplication by the fixed coefficients and finite addition preserve
the derivative class, so the resulting sum satisfies
\eqref{eq:profile-derivatives} with a constant depending on the choices.
It is compactly supported and real, with signs allowed, so it is an
averaging profile.
\end{proof}

\begin{proof}[Proof of Theorem~\ref{thm:main}]
First suppose \(|f_j|\leq1\), and write
\[
 F_n(x)=\prod_{j=1}^3 f_j(T^{jn}x),\qquad
 p=\prod_{j=1}^3\int_Xf_j\,d\mu,\qquad
 A_N(x)=\frac1N\sum_{n=1}^NF_n(x).
\]
For each fixed averaging profile \(\varphi\),
Proposition~\ref{prop:orbit-oscillation},
\eqref{eq:smooth-uniform-bound}, and Lemma~\ref{lem:ergodic-cauchy}
give a measurable conull set \(E_\varphi\) on which \(V_k^\varphi\) has a
limit \(V_\infty^\varphi\).  Extend that limit by zero off \(E_\varphi\);
it is measurable.  The mixing hypothesis and
Lemma~\ref{lem:ergodic-smooth-norm}, with slopes \(1,2,3\) and \(a=2^k\),
give \(\|V_k^\varphi-p\|_2\to0\).  Fatou's lemma on the convergence set
then yields
\[
 \int_{E_\varphi}|V_\infty^\varphi-p|^2\,d\mu
 \leq\liminf_{k\to\infty}\|V_k^\varphi-p\|_2^2=0,
\]
so \(V_k^\varphi\to p\) on a measurable conull set for each fixed profile.

For \(c\in[1,2]\), \(a=2^k\), and \(N_k=\lfloor ca\rfloor\), the
indicator convention in \(h_c\) gives
\[
 \sum_{n\in\mathbb Z}a^{-1}h_c(n/a)F_n(x)
   =\frac1{ca}\sum_{n=1}^{N_k}F_n(x)
   =\frac{N_k}{ca}A_{N_k}(x).
\]
Since \(|A_N|\leq1\) and \(0\leq ca-N_k<1\), for every averaging
profile \(\varphi\) and every \(x\),
\begin{equation}
 |A_{N_k}(x)-V_k^\varphi(x)|
 \leq\frac1{c2^k}
       +2^{-k}\sum_{n\in\mathbb Z}
          |h_c(n/2^k)-\varphi(n/2^k)|.
 \label{eq:ergodic-approximation}
\end{equation}
The function \(|h_c-\varphi|\) is bounded and compactly supported, with
possible discontinuities only at \(0,c\).  It is Riemann integrable, and
the sum on the right, including its factor \(2^{-k}\), tends to
\(\|h_c-\varphi\|_1\).

For each rational \(c\in[1,2]\) and integer \(h\geq1\), choose once an
averaging profile \(\varphi_{c,h}\) with
\(\|h_c-\varphi_{c,h}\|_1<2^{-h}\), using
Lemma~\ref{lem:ergodic-moments}.  Intersect their measurable conull
convergence sets over this countable family, and call the resulting
measurable conull set \(E\).  For \(x\in E\), every such \(c\), and
every \(h\), \eqref{eq:ergodic-approximation} gives
\[
 \limsup_{k\to\infty}|A_{\lfloor c2^k\rfloor}(x)-p|
 \leq\|h_c-\varphi_{c,h}\|_1<2^{-h}.
\]
Letting \(h\to\infty\), we obtain
\begin{equation}
 A_{\lfloor c2^k\rfloor}(x)\longrightarrow p
 \quad\text{for every }x\in E
       \text{ and every rational }c\in[1,2].
 \label{eq:ergodic-rational-lengths}
\end{equation}

Fix an integer \(s\geq1\) and the rational grid
\(c_i=1+i/s\), \(0\leq i\leq s\).  For \(2^k\leq N<2^{k+1}\), choose
the largest \(i\) for which \(n=\lfloor c_i2^k\rfloor\leq N\).
Then \(i<s\), \(n\geq2^k\), and
\[
 0\leq N-n
   <\lfloor c_{i+1}2^k\rfloor-\lfloor c_i2^k\rfloor
   \leq 2^k/s+1.
\]
For any \(1\leq n\leq N\), boundedness of the summands gives
\[
 |A_N-A_n|
 \leq\left(\frac1n-\frac1N\right)n+\frac{N-n}{N}
 =\frac{2(N-n)}{N}.
\]
For the selected \(n\), this is at most \(2/s+2/2^k\).
On \(E\), the maximum of
\(|A_{\lfloor c_i2^k\rfloor}(x)-p|\) over the finite grid tends to zero
by \eqref{eq:ergodic-rational-lengths}.  Consequently
\[
 \limsup_{N\to\infty}|A_N(x)-p|\leq2/s\qquad(x\in E).
\]
This holds for every \(s\geq1\), so the limsup is zero.

For general bounded measurable complex inputs, choose finite
\(M_j\geq\max(1,\sup_x|f_j(x)|)\), apply the result to \(f_j/M_j\), and
multiply by \(M_1M_2M_3\).  The conull set may depend on the original
triple.  Every exceptional set used in the proof is measurable and every
intersection is countable, so the proof requires no completeness or
standardness assumption on the probability space.
\end{proof}
